\documentclass{subfiles}
\begin{document}\label{sec:huffmanCodes}
    The first method to build optimal codes was discovered by D. A. Huffman,
        as solution to a problem posed by his lecturer Robert Fano\footnotemark.
        

    Before we discuss such a method, let us observe some properties of optimal codes:
        \begin{enumerate}
            \item \(\forall i,j, p_{i} \ge p_{j} \implies \abs{c_{i}} \le \abs{c_{j}}\).
                We could otherwise swap the codeword and obtain a code with a lower ACL.

            \item The least probable symbols have the same length.
                It's easely noted that if not, we could find a code with a lower ACL.

            \item The longest codewords differ on the last symbols.
                It follows directly from the previous property.
        \end{enumerate}

    The procedure itself is pretty straightforward; 
        it builds a \(d\)-ary tree, \(d\) being the size of the code alphabet,
        label each brach with a different symbols of the code alphabet and 
        assigns the codewords by traversing the tree.
        Practically speaking, we can summarize the algorithm as follow:
    \begin{enumerate}
        \item If not available, compute the symbols probabilities by scanning the text.
        \item Sort the symbols according to their probabilities, non-increseingly.
        \item Make sure that the number of source symbols \(n\) is \(d + \alpha(d - 1)\),
            for some integer \(\alpha\).
            If needed add \emph{dummy} symbols with probability 0.
        \item Merge the least \(d\) probable symbols in to a node until a root is formed.
        \item Assign the codeword by traversing the tree.
    \end{enumerate}

    Is it obvious that encoding a text with Huffman has a significant overhead
        depending on the code alphabet. We must infact transmit, in addition to the 
        bit stream, also the tree. Decoding the code is straightforward as well: 
        we simply traverse the tree and decode each symbol. 

    In terms of complexity Huffman encoding/decoing takes \(\bigO{n \log n}\),
        where \(n\) is the number of source symbols. 

    In addition to the tree overhead, we have to consider the possibility in 
        which the tree gets corrupted. In the next few sections we consider two
        variats of the plain Huffman encoding that try to fix these issues.

    \footnotetext{At that time,
        Fano and Shannon were looking for a technique to build optimal codes.}
    \clearpage
    \begin{example*}
        Consider the source \(S = \set{a, b, c, d, e, f, g}\), 
            where \(\Prob{a} = 0.3, \Prob{b} = 0.125\), \(\Prob{c} = 0.125, 
            \Prob{d} = 0.05, \Prob{e} = 0.1, \Prob{f} = 0.1, \Prob{g} = 0.2\).
            Provide a ternary Huffman code for \(S\).

        To begin with, let us sort the source symbols; we have: \(a \succeq g \succeq b
            \succeq c \succeq e \succeq f \succeq d\). Since the number of symbols
            is of the form \(d + \alpha(d - 1) = 3 + 2*2\), we can move on with merging the nodes.
            Thus, we consider the three least probable symbols and we merge them;
            in our case the symbols \(d, e, f\). 
            Iterating the procedure we get the tree in \Cref{fig:exampleHuffman}.

            \subfile{../../TikZ/Figure 2 - Example of Huffman encoding}

        Therefore, we get the code \(\C* = \set{0, 10, 11, 12, 20, 21, 22}\).

        Assume now receiving the stream \(22 21 22 11 11 0 10 20\); 
            knowing that the tree used is the one above, what is the encoded string?
            By trivially traversing the tree, we can easely decode the stream as 
            \(dfdbbage\).
    \end{example*}

    \subsection{Canonical Huffman encoding}
    \subfile{../sub/2.1 - Canonical Huffman encoding.tex}

    \subsection{Adaptive Huffman encoding}
    \subfile{../sub/2.2 - Adaptive Huffman encoding.tex}

    \nocite{huffman1952}
\end{document}
